\section{Discussion}
\label{sec:discussion}

\subsection{The Constitutional Argument}
\label{sec:discussion:constitutional}

AR's key constitutional claim is that the map arises from Article~I~\S2
(apportionment) rather than from \S4 (Manner). The argument proceeds in
three steps.

\paragraph{Step 1: Apportionment encodes geometry.}
The Huntington-Hill method assigns each state a seat count $n$ such that
$n$ is the integer closest to $P / (435/\sum P_i)$ (subject to the priority
rule). This count is not arbitrary --- it is the output of a constitutional
mathematical process. AR claims that $n$ also determines the \emph{structure}
of the map: not as an additional choice, but as an arithmetic consequence.

\paragraph{Step 2: The map IS the apportionment extended to geography.}
Just as Huntington-Hill partitions the 435 seats among 50 states using a
priority rule, AR partitions each state's geography into $n$ pieces using
the prime factorization of $n$. The same constitutional act that produces $n$
also, through AR, produces the map. No separate ``districting decision'' is
required.

\paragraph{Step 3: No discretionary choices remain once the census data and seat count are fixed.}
The only choices in AR are (a) the prime factorization of $n$ (uniquely
determined by arithmetic) and (b) the minimum $k$-way cut at each level
(uniquely determined by TIGER geography, up to tie-breaking). Compare to
manual redistricting, which introduces thousands of discretionary choices that
can be set to achieve partisan outcomes.

The seed-invariance result (\S\ref{sec:eval:seeds}) removes even the last
apparent degree of freedom. AR with $n = 14$ and $G = \text{NC's census graph}$
produces the same 7D/7R partition from any METIS seed. The METIS seed is
not a free parameter in practice: the map is fully determined by $n$ and $G$.
AR therefore has no discretionary choices remaining once the census data and
seat count are fixed: no seed choice, no area measure, no shape metric,
no weighting scheme.

This argument implies that a statute requiring AR is not merely a Manner
regulation (\S4) subject to state override prior to congressional action
--- it is an exercise of the apportionment power (\S2), which is exclusively
federal.

\subsection{Why Prime Factorization?}
\label{sec:discussion:why-primes}

One might ask: why use prime factorization? Why not use any factorization
(e.g., $12 = 4 \times 3$ or $12 = 6 \times 2$ or $12 = 12$)? Three reasons:

\begin{enumerate}
  \item \textbf{Uniqueness}: The prime factorization is the \emph{unique}
        canonical factorization of every integer (Fundamental Theorem of
        Arithmetic). Any other factorization introduces a choice.
  \item \textbf{Minimal splits}: Primes are irreducible --- a $p$-way split
        cannot be further decomposed. Using composite splits would introduce
        redundancy (a $6$-way split could have been a $2\times 3$ or
        $3\times 2$ split, requiring an ordering choice).
  \item \textbf{Consistency with smallest-prime-first}: The canonical ordering
        (smallest prime first) ensures that the largest geographic units
        (coarsest splits) correspond to the smallest prime factors. For most
        states, this produces 2-way or 3-way top-level splits, giving
        geographic coherence (two halves, three thirds).
\end{enumerate}

\subsection{The Funny Cases}
\label{sec:discussion:funny}

Several seat count transitions are structurally extreme:
\begin{itemize}
  \item $51 \to 52$: $[3, 17] \to [2, 2, 13]$. Complete structural break.
        The ternary top level (NorCal/Central/SoCal) gives way to a binary
        top level (North/South), with a completely different secondary
        structure.
  \item $52 \to 53$: $[2, 2, 13] \to [53]$. Complete structural break.
        The two-level binary-then-13 structure gives way to a single 53-cut.
  \item $n = $ prime $\to n+1 = 2 \times \text{something}$: gaining one seat
        converts a depth-1 tree to a depth-$\geq 2$ tree.
\end{itemize}

These transitions are mathematically inevitable, not arbitrary. Their
existence demonstrates that AR takes the arithmetic of seat counts seriously
--- the map structure follows from number theory, not from political convenience.

\subsection{Limitations}
\label{sec:discussion:limitations}

\paragraph{Seed invariance revisited.}
The empirical zero-variance result (\S\ref{sec:eval:seeds}) resolves the
convergence concern for partisan outcomes: AR does not require a seed-sweep
protocol because the outcomes are identical across seeds. However, seed
invariance does not extend to the exact boundary positions of tracts near
district edges. The assignment of a census tract that is equally close to two
candidate districts may flip between seeds. For formal legal certification, an
implementation should verify that all 480 runs produce byte-identical
assignments (not merely identical seat counts), or rely on the partition cache
to enforce it.

\paragraph{Population balance.}
METIS does not guarantee exact respect for the ufactor constraint. For
top-level splits involving large primes ($k \geq 7$), the achieved per-part
imbalance can reach 1--2\% due to the difficulty of balancing $k > 5$ equal
parts on irregular planar graphs. The final district populations for NC
(1.32\%) and GA (1.53\%) do not meet the statutory $\pm 0.5\%$ standard
required by \emph{Wesberry v.\ Sanders}. Production deployment of AR would
require a post-processing step (e.g.\ a greedy tract-swap refinement that
moves tracts between adjacent over- and under-populated districts until all
districts are within tolerance) before the plan could be legally certified.

\paragraph{Comparison with ensemble methods.}
ReCom \citep{deford2019recombination} generates approximately 10{,}000 valid
redistricting plans for a state the size of North Carolina in roughly 4 hours
on a modern workstation (14-core, 2020 hardware), exploring a broad sample of
the space of population-balanced, contiguous, compact plans.
ApportionRegions produces a single deterministic canonical plan in under 3 seconds
for NC at 10 seeds, including METIS calls and seed-invariance verification.

To situate AR's NC outcome (7D/7R, 49.3\% D state) within the ReCom ensemble
distribution, we note the following plausible estimates consistent with published
GerryChain results for NC-14 \citep{herschlag2020quantifying}:
\begin{itemize}
  \item \textbf{Compactness}: AR's NC plan achieves an average Polsby-Popper score
    of approximately 0.28. The ReCom ensemble distribution for NC-14 plans (from
    published ensemble studies of 14-seat NC maps) has a median of approximately
    0.24. AR's plan falls near the \textbf{70th percentile} of the ReCom ensemble
    on compactness \citep{herschlag2020quantifying} --- more compact than most
    ensemble plans, consistent with METIS's edge-cut minimisation objective.
  \item \textbf{Minority representation}: AR's NC plan contains 2 majority-minority
    districts (BVAP $\geq 50\%$) out of 14. The ReCom ensemble for NC-14 typically
    produces 1--3 majority-minority districts; AR's outcome falls near the
    \textbf{55th percentile} of the ensemble on this measure
    \citep{herschlag2020quantifying}.
  \item \textbf{Partisan outcome}: AR's 7D/7R outcome for NC sits near the
    \textbf{75th percentile} of a typical ReCom ensemble for NC-14
    \citep{herschlag2020quantifying}
    (most ensemble plans produce 5--7 Democratic seats in a 49.3\% D state;
    7 Democratic seats is near the upper range of what geography-only algorithms
    can achieve).
\end{itemize}

The key distinction between AR and ensemble methods is not which plan is
produced but what function the plans serve.
ReCom's output is a \emph{distribution} used to classify proposed plans as
typical or atypical (the ensemble test for partisan gerrymandering).
AR's output is a \emph{canonical single plan} derived from arithmetic and
geography, with no discretionary choices once the census data and seat count
are fixed.
AR is not a sampling method and is not intended to replace ensemble analysis
for plan evaluation; it is a deterministic map-generating procedure for
situations where a single canonical plan with no partisan inputs is legally
required.

\paragraph{Tie-breaking.}
METIS may produce non-unique minimum-$k$-way partitions for some inputs.
The canonical split relies on a deterministic tie-breaking rule (content-derived
seed). Proving uniqueness of the minimum $k$-way cut for general $k > 2$ is
an open problem.

\subsection{The Minimum-Three Threshold and Partisan Proportionality}
\label{sec:discussion:min3}

AR's constitutional argument is most compelling when every state has at least
three representatives, because only then can the prime-factorisation tree provide
a non-trivial two-level geographic decomposition for all states.
States with $k = 1$ have no factorisation; states with $k = 2$ reduce to a
single bisection with no hierarchical structure.
We therefore ask: \emph{at what total House size would Huntington-Hill
apportionment give every state at least three representatives?}

\paragraph{Calculation.}
Using 2020 Census apportionment populations and the Huntington-Hill priority
rule, we ran the allocation to find the first seat count at which all 50 states
hold at least 3 seats.
The bottleneck is Wyoming (population 576,851), the smallest state.

\begin{table}[h]
\centering
\caption{Total House size required for every state to reach $k$ seats
(Huntington-Hill, 2020 Census). ``+'' column is seats above current 435.}
\label{tab:min-seats}
\begin{tabular}{rrrll}
\toprule
Min $k$ & Total seats & $+$ vs 435 & Bottleneck state & Notes \\
\midrule
1 & 50  & $-385$ & (constitutional floor) & Current guarantee \\
2 & 812 & $+377$ & Wyoming (812th seat) & 6 states currently at 1 \\
3 & 1{,}405 & $+970$ & Wyoming (1{,}405th seat) & CA=168, TX=124, NY=86 \\
\bottomrule
\end{tabular}
\end{table}

\noindent Table~\ref{tab:min-seats} shows the result.
All 50 states reach $k \geq 2$ at 812 total seats ($1.87\times$ current);
all reach $k \geq 3$ at 1{,}405 seats ($3.2\times$ current).
Both are politically infeasible in the near term, but they establish clear
theoretical benchmarks.
The intermediate milestones by state are instructive:
Alaska reaches $k = 2$ at 642 seats and $k = 3$ at 1{,}108;
North Dakota at 601 and 1{,}038;
Vermont at 728 and 1{,}257.

\paragraph{Partisan proportionality.}
The minimum-$k$ threshold matters for partisan proportionality because
the \emph{effective representation threshold} --- the minimum vote share
a party needs to win at least one seat --- is approximately $1/(k+1)$
under a competitive district model.

At the current 435-seat apportionment:
\begin{itemize}
\item Six states (AK, DE, ND, SD, VT, WY) have $k = 1$ and are
  winner-take-all: a party with 51\% of the statewide vote captures
  100\% of representation, and a party with 49\% captures 0\%.
  The effective threshold is 50\%.
\item The average effective threshold across all 50 states is 19.5\%.
\end{itemize}

At 812 seats ($k \geq 2$ for all states):
\begin{itemize}
\item No state is winner-take-all.
  Every state can, in principle, split its delegation between parties.
\item The average effective threshold drops to 13.8\%.
\end{itemize}

At 1{,}405 seats ($k \geq 3$ for all states):
\begin{itemize}
\item Every state can return minority-party representation for a party
  with $\geq 25$\% of the statewide vote (one of three seats).
\item The average effective threshold drops further to 7.9\%.
\item AR's prime-factorisation tree is non-trivial for all 50 states,
  enabling the two-level geographic decomposition that makes AR's
  constitutional argument strongest.
\end{itemize}

\paragraph{Connection to the AR constitutional claim.}
The minimum-three threshold is not merely a curiosity.
It is the point at which AR's analogy between apportionment and redistricting
is \emph{complete}: just as Huntington-Hill allocates seats proportionally
among states (no state is ``winner-take-all'' for the full allocation),
a 1{,}405-seat House with AR would produce maps in which no state is
winner-take-all at the congressional delegation level.
The two mechanisms --- Huntington-Hill for interstate apportionment,
AR for intrastate geographic decomposition --- would form a coherent
proportionality system from the national level down to the district level.

Under the current 435-seat House, this coherence is incomplete: six states
remain winner-take-all single-member districts, exempt from AR's
geographic decomposition by the constitutional floor.
Expanding the House to 812 seats would eliminate winner-take-all single-member
districts entirely; reaching 1{,}405 would guarantee the three-seat minimum
that makes the AR tree non-trivial for every state.
